\begin{textsample}{POS Dim 1 – vem\_ed\_07 – Score -1.00 – vem\_ed\_07/t025.txt}  \label{ex:f1_pos_186}
QUEBRA-GELO QUEBRA-GELO A equipe dos Projetos Colaborativos Internacionais (PCIs) da Unidade de Ensino Superior (Cesu) do Centro Paula Souza vem buscando diversificar parcerias com instituições de ensino \textbf{superior} (IES) em países latinoamericanos.
Entre elas, está a universidade colombiana Uniminuto, com a qual \textbf{foram} desenvolvidos quatro projetos entre agosto de 2020 e maio de 2021, nas Fatecs Barueri, Piracicaba, Guaratinguetá e Itaquaquecetuba.
Os temas abordados envolveram \textbf{competências} emocionais, aspectos de cultura e educação, publicidade digital e marketing nos dois países.
Essas colaborações ocorreram em língua espanhola, mas a Uniminuto oferece cursos de Letras / Língua Portuguesa, e demonstrou interesse em realizar projetos no nosso \textbf{idioma}.
A primeira iniciativa para atender essa demanda \textbf{foi} a realização do Colóquio de Literatura Brasileira, oferecido aos estudantes da Uniminuto pelas Fatecs Itaquaquecetuba e Guaratinguetá, em 14 de maio.
Nas páginas 2 e 3, você fica \textbf{sabendo} mais sobre as colaborações entre Fatecs e Uniminuto.
Na seção “Quem \textbf{é} Quem”, temos o depoimento de James Goodwin, diretor de iniciativas da BridgeValley Community and Technical College (EUA).
Em 2019 e 2020, a instituição desenvolveu um projeto sobre equipamentos médico-hospitalares com alunos de Sistemas Biomédicos da Fatec Bauru.
E, desde o ano passado, realiza um projeto sobre física das ondas com estudantes de Produção Fonográfica da Fatec Tatuí.
Em “Boas Práticas”, apresentamos o relato de Camila Kami, professora de inglês das Fatecs Bauru e Garça.
Experiente, Camila atua com PCIs desde o segundo semestre de 2019.
Boa leitura!

% matched lemmas: competência, idioma, saber, ser, superior
\end{textsample}
